\section*{Part VIII: Asset Management and the Long View}
\addcontentsline{toc}{section}{Part VIII: Asset Management and the Long View}


\subsection{61. A completed building creates a continuing liability}

A facility requires periodic repairs, replacement of HVAC, networking, transformers, interiors, fire systems, roofs, roads and lab equipment. Initial capex is only the beginning. Fund a \textbf{capital renewal reserve} rather than assuming all tuition after opening can support new expansion. Deferred maintenance can become a severe safety and reputational failure years later.

Track asset condition, lifecycle replacement, downtime, energy/water efficiency, regulatory compliance and utilization. Operations and capex planning should speak to each other.


\subsection{62. Total cost of ownership}

For each building or site calculate acquisition, development, finance, taxes, insurance, maintenance, utilities, people, compliance, renewals, future demolition or remediation, and credible residual/exit value. Use present-value scenarios where useful. A cheap low-quality structure can have higher lifetime cost than a resilient design.

For specialized laboratories include instrument calibration, safety supplies, waste management and obsolete equipment disposal. For student housing include staffing, round-the-clock response, health and welfare costs.


\subsection{63. Campus utilization: avoid vanity denominators}

Report scheduled use, actual attendance, research output, room-hour utilization and \textbf{cost per relevant outcome}. A classroom used 20 hours per week may still be justified if quality and demand require it; a building used 90 hours per week may be a poor investment if the activity produces little learning. Report peak constraints separately from average occupancy to avoid overbooking laboratories.

Facilities must be assessed against mission outputs and full cost, not just carpet area.


\subsection{64. The property portfolio and future mergers}

An institute might accumulate multiple small facilities over time and later discover they hinder formation of a single academic community. Before each lease or purchase, ask whether it fits a future hub-and-spoke strategy, how it could be exited, and whether it creates stranded staff and obligations. Do not let historical sunk costs dictate the optimal future campus footprint.

If an acquired institution brings property, conduct the same title, compliance and asset-condition review as for any new purchase.


\subsection{65. Insurance, disaster recovery and business continuity}

Evaluate property, construction all-risk, 3rd-party liability, equipment, business interruption, employee and student risks and specialized research exposures, with counsel and brokers. Maintain physical emergency plans, alternate teaching arrangements, backups for research data, asset inventories and critical vendor contacts.

A credible institution can continue essential teaching after a flood, fire, power outage or major contractor failure. Insurance may reduce loss severity but does not replace duty of care.
